% Motivation and learning outcomes.
% The outcomes are the outcomes of Day 9 in the syllabus.
\begin{frame}{Why this day matters}
  \begin{itemize}
    \item Days 3 to 8 gave many numbers: a latency, a model size, an accuracy,
          a frame rate. A number is useful only when a second person can
          measure it again and get the same value.
    \item One timed inference is not a measurement. The first inferences are
          slower than the later ones, the latency changes from run to run, and
          a mean hides the slow inferences. A fair benchmark has a warm-up,
          many repetitions, fixed inputs, and percentiles.
    \item A battery device has an energy budget. The energy of one inference
          and the duty cycle give the battery life. A hot processor reduces
          its clock frequency, so a short run gives a latency that a long run
          does not keep.
    \item In the lab you write the benchmark protocol first. Then you measure
          keyword spotting and image classification on the XIAOML Kit, and
          image classification and object detection on the Raspberry Pi 5. You
          write one report for the two boards.
  \end{itemize}
\end{frame}

\begin{frame}{Learning outcomes}
  After this day you can:
  \begin{itemize}
    \item Define latency percentiles, throughput, peak memory, and energy per
          inference.
    \item Design a fair benchmark with warm-up, repetitions, and fixed inputs.
    \item Measure one task on two boards and explain the differences.
    \item Explain thermal throttling and its effect on results.
  \end{itemize}
\end{frame}
